\printedpage{69}
% head: A. Grothendieck : Dérivateurs, Ch. XV, Édition des Universités de Montpellier II et Paris VII
% foot: Version 5 novembre 2010
where
\[
\pr_2^{I_\alpha}\ :\ I_\alpha\times X \longrightarrow X\ .
\]

\msp{page 128}Let us now consider the diagram (solid lines)
\[
\begin{tikzcd}[column sep=small]
& X\amalg X \arrow[dl, "{(\partial_0^X,\partial_1^X)}"'] \arrow[dr, "{(\partial_0^{\prime X},\partial_1^{\prime X})}"] & \\
I\times X \arrow[rr, "\varphi\times\id_X\,=\,\varphi_X"] \arrow[dr, "h"'] & & J\times X \arrow[dl, dotted, "h'"] \\
& Y\ , &
\end{tikzcd}
\]
where $\partial_0^{\prime X}$, $\partial_1^{\prime X}$ are defined as $\partial_0^X$, $\partial_1^X$, but using $j_0,j_1\in\Ob J$ instead of $i_0,i_1\in\Ob I$. If there existed in Cat an arrow $h'$ (dotted) such that $h'\circ\varphi_X=h$, it would follow that $h'$ is a homotopy of $f_0$ and $f_1$ along $(J,j_0,j_1)$. Since $J$ is 0-connected, it would follow at once that $f_0$, $f_1$ are homotopic in the usual sense in Cat, hence that $f_0$ and $f_1$ have the same image in $\Hot_W$. But if one cannot find an arrow $h'$ in Cat, at least one can find it in $\Hot_W$, since $\varphi_X\in W$, hence an isomorphism in $\Hot_W$ (and $h'$ in $\Hom_W(J\times X,Y)$ will even be uniquely determined). Thus in $\Hot_W$, we have
\[
f_0\ =\ h'\partial_0^{\prime X}\ ,\qquad\quad f_1\ =\ h'\partial_1^{\prime X}\ .
\]

\msp{page 129}To prove that
\[
f_0\ =\ f_1\qquad\text{in}\quad\Hom_W(X,Y)\ ,
\]
it suffices to see that
\[
\partial_0^{\prime X}\ =\ \partial_1^{\prime X}\qquad\text{in}\quad\Hom_W(X,J\times X)\ .
\]
Now in fact $\partial_0^{\prime X}$ and $\partial_1^{\prime X}$ are homotopic, for $e\underset{j_1}{\overset{j_0}{\rightrightarrows}}J$ are homotopic ($j_0$ and $j_1$ belonging to the same connected component of $J$), hence also their product with $\id_X$.

One sees therefore that for any two objects $X$, $Y$ of Cat, all the arrows from $X$ to $Y$ are equal in Hot. The argument given even proves that all the arrows \textit{in} $\Hot_W$ from $X$ to $Y$ are equal. This means that the category $\Hot_W^*$ (where for any two objects $X$, $Y$, $\Hom_{\Cat}(X,Y)$, and a fortiori $\Hom_W(X,Y)$ is $\neq\emptyset$) is equivalent to the one-point category. As for $\Hot_W$ itself, whose set of objects is $\{\emptyset\}\amalg\Ob\Hot_W^*$, with $\emptyset$ an initial object, it will therefore be equivalent to the one-point category, or else to $\boldsymbol{\Delta}_1$, according as $\emptyset\longrightarrow e$ is an isomorphism in $\Hot_W$, or not.

\msp{page 130}\textbf{Comment 9.8.} If the axioms Loc\,(1, 2, 3, 5)~\footnote[52]{In fact, this comment holds even in the absence of the assumption Loc\,(5), \textit{cf.} 9.9 below.} are satisfied, then the connectedness axiom Loc\,(8) \textit{is equivalent} to each of the following:

\printedpage{70}
% head: A. Grothendieck : Dérivateurs, Ch. XV, Édition des Universités de Montpellier II et Paris VII
% foot: Version 5 novembre 2010
\begin{enumerate}[leftmargin=*]
\item[\textit{a})] $\Hot_W$ is equivalent neither to $e$, nor to $\boldsymbol{\Delta}_1$.
\item[\textit{b})] $\Hot_W$ has infinitely many isomorphism classes.
\item[\textit{c})] $f_0:X_0\longrightarrow Y_0$ (arrow given in advance in Cat, such that $\pi_0(f_0)$ is not an isomorphism, and $X_0\neq\emptyset$) is \textit{not} an isomorphism in $\Hot_W$.
\item[\textit{d})] $X_0$ and $Y_0$ are not isomorphic in $\Hot_W$ (where $X_0$, $Y_0$ are two objects given in advance in Cat, with $X_0$, $Y_0$ nonempty and $\card\pi_0(X_0)\neq\card\pi_0(Y_0)$).
\end{enumerate}

% marginal mark: a double vertical rule in the left margin beside the first two lines of this paragraph
\marginbar One can therefore say that the connectedness axiom plays the role of a \textit{nontriviality axiom} for the localized category $\Hot_W$ of Cat. It will be satisfied automatically for all the basic localizers $W$ that we will be led to construct or to introduce. But we will practically never have to invoke it as an assumption, with a view to certain conclusions.

I finally realize that in proposition 9.7 the assumption Loc\,(5) is unnecessary:

\textbf{Proposition 9.9.}~\footnote[53]{\textit{Cf.} better statement 9.12 (p. 133 bis).} \textit{Let $W$ be a basic localizer in} Cat (i.e. \textit{we have} Loc\,(1, 2, 3) \textit{and nothing more}), \textit{plus the condition that the} sum \textit{of $f_1,f_2\in W$ is in $W$}~\footnote[54]{This condition Loc\,(5\,a) holds in particular if $W$ also satisfies one of the two axioms Loc\,(4) (\textit{cf.} 9.1), or Loc\,(6\,c$'$) (\textit{cf.} 9.3).}. \textit{If the connectedness axiom is not satisfied, then} $\Hot_W$ \textit{is equivalent to the one-point category} (\textit{if $\emptyset\longrightarrow e$ is an isomorphism in} $\Hot_W$~\footnote[55]{This also means that there exists $Z\neq\emptyset$ in Cat, such that $\emptyset\longrightarrow Z$ is an isomorphism in $\Hot_W$.}) \textit{or to $\boldsymbol{\Delta}_1$} (\textit{in the opposite case}), \textit{and} $\Hot_W^*$ ($\overset{\text{def}}{=}\Hot_W\backslash\{\emptyset\}$) \textit{is equivalent to the one-point category.}

\msp{page 131}The condition on sums ensures that in $\Hot_W$ the sum of two objects exists, and that $\Cat\longrightarrow\Hot_W$ commutes with sums of two objects. This said, let
\[
f\ :\ X\longrightarrow Y \tag*{$(*)$} % the label (*) is printed at the left margin
\]
be an arrow such that $f\in W$, but $\pi_0(f)$ not surjective. Let us first suppose $X\neq\emptyset$ (the least obvious case). If $Y_0$ is a nonempty summand of $Y$ such that $f$ factors through $Y_1=Y\backslash Y_0$, then $f$ is the sum of
\[
f_1\ :\ X\ =\ X_1\longrightarrow Y_1\qquad\quad\text{and}\qquad\quad f_0\ :\ X_0\ =\ \emptyset\longrightarrow Y_0
\]
in Cat, hence also in $\Hot_W$. Thus we have a sum $f\ =\ f_0\amalg f_1$ in $\Hot_W$ which is an isomorphism, hence for every $Z$ in $\Hot_W$,
\[
\Hom_W(Y_0,Z)\times\Hom_W(Y_1,Z)\ \xrightarrow{\ \sim\ }\ \underbrace{\Hom_W(X_0,Z)}_{=\,e}\times\Hom_W(X_1,Z)
\]
(\textbf{N.B.} $\emptyset$ is an initial object of $\Hot_W$, being so in Cat), \textit{i.e.} the composite
\[
Z(Y_0)\times Z(Y_1)\ \xrightarrow{\pr_2}\ Z(Y_1)\ \xrightarrow{Z(f_1)}\ Z(X_1)
\]
is bijective. Since $X\ =\ X_1\neq\emptyset$, hence also $Y_1\neq\emptyset$, and $Y_0\neq\emptyset$, the sets $Z(Y_0)$, $Z(Y_1)$, $Z(X_1)$ are $\neq\emptyset$ if $Z\neq\emptyset_{\Cat}$. The composite being bijective, the first arrow $\pr_2$ is injective, which implies
\printedpage{71}
% head: A. Grothendieck : Dérivateurs, Ch. XV, Édition des Universités de Montpellier II et Paris VII
% foot: Version 5 novembre 2010
\msp{page 132}that $Z(Y_0)$ is reduced to a point (the two factors $Z(Y_0)$, $Z(Y_1)$ being $\neq\emptyset$). This means that $Y_0$ is an initial object in $\Hot_W^*$. It follows in particular that $Y_0\amalg Y_0$ is also an initial object, and the canonical map $Y_0\amalg Y_0\longrightarrow Y_0$ is therefore an isomorphism in $\Hot_W$, whereas the induced map on the $\pi_0$ is not injective. (It is $\pi_0(Y_0)\amalg\pi_0(Y_0)\longrightarrow\ \pi_0(Y_0)$, with $\pi_0(Y_0)$ nonempty.) This case will be treated below.

If in $(*)$ $X$ is empty, $Y$ nonempty (since $\pi_0(f)$ not surjective), then $Y$ is an initial object of $\Hot_W$, and one sees again that $Y\amalg Y\longrightarrow Y$ is an isomorphism in $\Hot_W$, although not inducing an isomorphism of the $\pi_0$. Proposition 9.9 will therefore follow from the

\textbf{Lemma 9.10.} \textit{Let $W$ be a basic localizer and nothing more} (\textit{axioms} Loc\,(1, 2, 3)). \textit{Suppose there exists $f:I\longrightarrow J$ in} Cat \textit{such that} $\hot_W(f)$ \textit{is an isomorphism, but $\pi_0(f)$} is not injective~\footnote[56]{Or else such that $I$ is 0-connected, and $J$ not 0-connected (\textit{i.e.} $\card\pi_0(I)\ =\ 1$, $\card\pi_0(J)\ \geq\ 2$). This second case is treated by the same method essentially.}. \textit{Then} $\Hot_W^*$ \textit{is equivalent to the one-point category $e$, and} $\Hot_W$ \textit{is equivalent to $e$} (\textit{if $\emptyset\longrightarrow e$ is an isomorphism in} $\Hot_W$) \textit{or to $\boldsymbol{\Delta}_1$} (\textit{in the opposite case}).

\msp{page 133}The proof is the one given for 9.7: one proves that if $X$, $Y$ are two objects of Cat, \textit{nonempty}, then two arrows $f_0,f_1:X\rightrightarrows Y$ are equal in $\Hot_W$, by a homotopy argument, choosing two points $i_0$, $i_1$ of $I$ which are in two distinct connected components, and such that their images $j_0$, $j_1$ in $J$ are in the same component.

\textbf{Proposition 9.11.} \textit{Let $W$ be a basic localizer} (i.e. \textit{we have} Loc\,(1, 2, 3)). \textit{Then the following conditions are equivalent.}
\begin{enumerate}[leftmargin=*]
\item[a)] \textit{We have} Loc\,(5\,bis) (\textit{decomposition axiom, strong form}) \textit{and} Loc\,(8) (\textit{connectedness axiom}).
\item[b)] \textit{We have} Loc\,(5\,bis), \textit{and} $\Hot_W$ \textit{is not equivalent to $e$ or to $\boldsymbol{\Delta}_1$} (cf. \textit{comment} 9.8).
\item[c)] \textit{For every $f:X\longrightarrow Y$, we have $f\in W$ if and only if $\pi_0(f)$ is an isomorphism, and for every connected component $X_i$ of $X$, if $Y_i$ denotes the corresponding connected component of $Y$, the induced arrow $f_i:X_i\longrightarrow Y_i$ is in $W$.}
\end{enumerate}

The reason why I do not isolate this property as one of the axioms Loc\,(\ldots) is that it is \textit{not} a stability axiom, but an axiom of a mixed nature.

\msp{page 133 bis}\textbf{Proposition 9.12.} \textit{Let $W$ be a basic localizer} (\textit{nothing more}), i.e. \textit{satisfying} Loc\,(1, 2, 3). \textit{The following conditions are equivalent.}
\begin{enumerate}[leftmargin=*]
\item[a)] \textit{$W$ satisfies the connectedness axiom} Loc\,(8), i.e. $f\in W\Longrightarrow\pi_0(f)$ \textit{bijective.}
\item[b)] \textit{If $f,g:X\rightrightarrows Y$ in} Cat \textit{are such that} $\hot_W(f)=\hot_W(g)$, \textit{then} $\pi_0(f)=\pi_0(g)$.
\item[c)] \textit{There exist $X$, $Y$ in} Cat \textit{such that the map} $\hot_W^{X,Y}:\Hom_{\Cat}(X,Y)\longrightarrow\Hom_W(X,Y)$ \textit{is not constant.}
\item[c$'$)] \textit{The subcategory} $\Hot_W^*$ \textit{of} $\Hot_W$, \textit{consisting of the nonempty objects, is not equivalent to the one-point category.}
\item[c$''$)] \textit{The category} $\Hot_W$ \textit{is equivalent neither to the category $e$ nor to the category $\boldsymbol{\Delta}_1$.}
\item[d)] \textit{The set of isomorphism classes of objects of} $\Hot_W$ \textit{is infinite.}
\end{enumerate}

\printedpage{72}
% head: A. Grothendieck : Dérivateurs, Ch. XV, Édition des Universités de Montpellier II et Paris VII
% foot: Version 5 novembre 2010
\textsc{Proof}. We have $(a)\Longrightarrow(b)$, since $(a)$ is equivalent to saying that the functor $\pi_0:\Cat\longrightarrow\mathrm{Ens}$ factors through $\Hot_W$. Moreover the corresponding functor
\[
\Hot_W\ \longrightarrow\ (\mathrm{Ens})
\]
is essentially surjective -- there is an inverse up to isomorphism, $I\longmapsto\hot_W(I)$. This proves that $(a)$ implies $(d)$, and we have $(d)\Longrightarrow(c')$, $(d)\Longrightarrow(c'')$ tautologically. A more or less tautological argument shows moreover that if $\Hot_W^*$ is equivalent to $e$, then $\Hot_W$ is equivalent to $e$ or to $\boldsymbol{\Delta}_1$ (for $\emptyset$ is an initial object of $\Hot_W$) -- hence $(c'')\Longrightarrow(c')$. Thus we have more or less tautologically
\[
\begin{tikzcd}
(\mathrm{a}) \arrow[r, Rightarrow] \arrow[d, Rightarrow] & (\mathrm{d}) \arrow[r, Rightarrow] & (\mathrm{c}'') \arrow[r, Rightarrow] & (\mathrm{c}') \\
(\mathrm{b}) \arrow[rr, Rightarrow] & & (\mathrm{c}) \arrow[ru, Rightarrow] &
\end{tikzcd}
\]
where the implications $(b)\Longrightarrow(c)\Longrightarrow(c')$ are also trivial \texttt{[it is rather the implication (c$'$) $\Longrightarrow$ (c) that is trivial~; the implication (c) $\Longrightarrow$ (c$'$) is proved after the statement of corollary 9.13, before the proof of the latter]}. Thus
\msp{page 133 ter}everything comes down to proving that $(c)\Longrightarrow(a)$, or again:

\textbf{Corollary 9.13.} \textit{If the connectedness axiom is not satisfied, then if $X$, $Y$ are two objects of} Cat, \textit{all the arrows $X\longrightarrow Y$ in} Cat \textit{are equal in} $\Hot_W$, i.e. \textit{the map} $\Hom_{\Cat}(X,Y)\longrightarrow\Hom_W(X,Y)$ \textit{is constant.}

This therefore implies that all the objects of Cat$^*$, \textit{i.e.} all the nonempty objects of Cat, are isomorphic in $\Hot_W$, and one finds a canonical isomorphism between them, with transitivity. Let then $\mathcal{H}$ be the quotient category of Cat, having the same objects as Cat, and having exactly one arrow between two nonempty objects of Cat and exactly one arrow from $\emptyset$ to every other object of Cat (and no arrow from a nonempty object of Cat to $\emptyset$). It is therefore a category equivalent to $\boldsymbol{\Delta}_1$. We have seen that $\Cat\longrightarrow\Hot_W$ factors as
\[
\Cat\ \xrightarrow{\text{can.}}\ \mathcal{H}\ \longrightarrow\ \Hot_W\ .
\]
If $\Cat\longrightarrow\mathcal{H}$ transforms every arrow belonging to $W$ into an isomorphism, then it is immediate that $\mathcal{H}$ is a solution to the universal problem defining $\Hot_W$, hence that $\mathcal{H}\longrightarrow\Hot_W$ is an isomorphism. This also means that $W$ does not contain an arrow $\emptyset\longrightarrow X$ with $X$ nonempty (in which case $W$ contains all arrows). Thus in this case $\Hot_W\simeq\boldsymbol{\Delta}_1$. In the opposite case, let $\mathcal{H}'$ be the category obtained from $\mathcal{H}$ by inverting the arrows $\emptyset\longrightarrow X$, \textit{i.e.} the 1-connected groupoid on the set of objects $\Ob\Cat$. It is no longer a quotient category of Cat, but we still have a factorization
\[
\Cat\ \xrightarrow{\text{epi}}\ \mathcal{H}\ \ \xrightarrow{\text{loc.}}\ \mathcal{H}'\ \longrightarrow\ \Hot_W\ ,
\]
\printedpage{73}
% head: A. Grothendieck : Dérivateurs, Ch. XV, Édition des Universités de Montpellier II et Paris VII
% foot: Version 5 novembre 2010
\msp{page 133 quater}and this time it is $\mathcal{H}'$ that is a solution of the universal problem ($\Cat\longrightarrow\mathcal{H}'$ transforming arrows of $W$ into isomorphisms). Thus in this case, $\mathcal{H}'\longrightarrow\Hot_W$ is an isomorphism, and $\Hot_W$ is equivalent to $e$. In both cases, $\Hot_W^*$ is equivalent to $e$.

Thus, proposition 9.12 follows from corollary 9.13, which it is now a matter of proving.

Let $f:I\longrightarrow J$ in Cat be such that $\pi_0(f)$ is not injective. We have seen in the proof of 9.7 how one deduces from this (by choosing two elements $i_0$, $i_1$ in two distinct connected components of $I$, such that $j_0=f(i_0)$ and $j_1=f(i_1)$ are in the same connected component of $J$) that for $X$, $Y$ in Cat, two arrows $f_0,f_1:X\rightrightarrows Y$ are always equal in $\Hot_W$ (\textit{cf.} lemma 9.10).

Let now
\[
f\ :\ I\longrightarrow J
\]
in Cat be such that $\pi_0(f)$ is not surjective, hence $J=J_0\amalg J_1$, with $f$ factoring through $J_0$, and $J_1\neq\emptyset$. Let us then consider the discrete category with two objects,
\[
\varepsilon\ =\ \{0,1\}\ ,
\]
and two arrows
\[
g_0,\ g_1\ :\ J\ \rightrightarrows\ \varepsilon\ ,
\]
$g_0$ being constant with value 0, and $g_1$ being equal to 0 on $J_0$, to 1 on $J_1$. We have
\[
g_0f\ =\ g_1f\ ,
\]
hence, since $f\in W$,
\[
g_0\ =\ g_1\qquad\quad\text{in}\quad\Hot_W\ .
\]
One concludes from this, by taking a point $j_1$ in $J_1$, and the corresponding inclusion $\alpha_{j_1}:e\longrightarrow J$, that
\msp{page 133 quinquies}
\[
g_0\alpha_{j_1}\ =\ g_1\alpha_{j_1}\ ,
\]
\textit{i.e.} the two sections of $\varepsilon$ over $e$ are equal in $\Hot_W$. Since the functor $\Cat\longrightarrow\Hot_W$ commutes with finite products \texttt{[this property is not established under the sole conditions Loc\,(1,2,3)~; it is proved in Pursuing Stacks, section 64, on replacing Loc\,(3) by the stronger axiom Loc\,(3\,bis)]}, it follows that for $X$ in Cat, the two arrows
\[
X\ \underset{\partial_1^X}{\overset{\partial_0^X}{\rightrightarrows}}\ X\times\varepsilon\ =\ X\amalg X
\]
are equal in $\Hot_W$. If then
\[
f_0,\ f_1\ :\ X\ \rightrightarrows\ Y
\]
are two arrows in Cat, they are of the form
\[
f_0\ =\ h\partial_0^X\ ,\qquad\quad f_1\ =\ h\partial_1^X\ ,
\]
with
\[
h\ =\ (f_0,f_1)\ :\ X\amalg X\ \longrightarrow\ X\ , % [sic] X\amalg X \longrightarrow X -> X\amalg X \longrightarrow Y, h\partial_i^X = f_i : X \to Y requires target Y
\]
and since $\partial_0^X=\partial_1^X$ in $\Hot_W$, one concludes $f_0=f_1$ in $\Hot_W$, \quad \textit{q.e.d.}